% ARCHETYPE: entry with header, subheader, and bullets.
% Two header lines -- bold title with right-aligned place, italic role with
% right-aligned dates -- then \cvlistgap, then the bullet list. The gap
% belongs after the header BLOCK, not after the first header line.
\newcommand{\sectionExperience}{
\section*{EXPERIENCE}
\textbf{Position or Project Title} \hfill City, Country\\
\textit{Role or Title} \hfill January 2024 -- December 2024
\cvlistgap
\begin{itemize}
    \item State what you did, the method or tool you used, and the outcome it produced. One sentence of roughly this length keeps the entry to a predictable height, which is what makes the rhythm visible.
\end{itemize}
\cventrygap
\textbf{Position or Project Title} \hfill City, Country\\
\textit{Role or Title} \hfill May 2023 -- December 2023
\cvlistgap
\begin{itemize}
    \item A second entry, so the entry-to-entry gap has something to separate. Note that nothing between these two entries creates space except \texttt{\textbackslash cventrygap}.
\end{itemize}
\cventrygap
\textbf{A Title Long Enough To Wrap Onto\\
    A Second Line} \hfill City, Country\\
\textit{Role or Title} \hfill May 2022 -- April 2023
\cvlistgap
\begin{itemize}
    \item A header may itself span two lines; the indented continuation above shows the convention. The list gap still measures from the end of the whole header block.
\end{itemize}
}
